\documentclass[11pt]{article}
\usepackage[margin=1in]{geometry}
\usepackage{amsmath,amssymb,mathtools}
\usepackage{booktabs,array}
\usepackage{xcolor}
\usepackage[hidelinks]{hyperref}
\usepackage{microtype}
\usepackage{fancyhdr}
\definecolor{navy}{RGB}{24,55,95}
\definecolor{green}{RGB}{35,112,75}
\pagestyle{fancy}
\setlength{\headheight}{14pt}
\fancyhf{}
\lhead{Dicyclic and generalized quaternion groups}
\rhead{Proof and verification}
\cfoot{\thepage}
\newcommand{\Dic}{\operatorname{Dic}}
\newcommand{\Sym}{\operatorname{Sym}}
\newcommand{\M}{\mathcal M}
\newcommand{\E}{\mathbb E}
\title{\color{navy}\textbf{Exact $\sigma$-MGFs for Dicyclic Groups}\\[2mm]
\large Proof, matrix construction, and finite verification}
\author{Computational proof companion}
\date{17 July 2026}

\begin{document}
\maketitle

\begin{abstract}
For the dicyclic group $\Dic_m$ of order $4m$, this note proves the proposed
generalized Molien formulas for every complex representation, specializes them
to the two-dimensional irreducibles $\rho_k$, and derives the stated exact
coefficient count.  It also checks direct sums and the regular representation.
The accompanying marimo notebook constructs the matrices themselves and
compares the rank of an explicit Reynolds projector with three independent
formula evaluations.  All representative cases passed to numerical error below
$2.6\times10^{-15}$.
\end{abstract}

\section{The general finite-group identity}
For a finite group $G$ acting on $V$ and a finite tuple
$\tau=(\tau_1,\ldots,\tau_s)$, put
\[
 W_\tau=\bigotimes_{b=1}^s\Sym^{\tau_b}V.
\]
If the eigenvalues of $\rho(g)$ are $\lambda_1,\ldots,\lambda_d$, then
\[
 \det(I-t\rho(g))^{-1}=\prod_{r=1}^d(1-t\lambda_r)^{-1}
 =\sum_{q\ge0}\chi_{\Sym^qV}(g)t^q.
\]
Consequently the coefficient of $m_\tau$ in the $\sigma$-MGF is
\[
 [m_\tau]\M_{G,V}
 =\frac1{|G|}\sum_{g\in G}\prod_b\chi_{\Sym^{\tau_b}V}(g)
 =\operatorname{tr}\!\left(\frac1{|G|}\sum_{g\in G}\rho_{W_\tau}(g)\right).
\]
The average in parentheses is the Reynolds projector onto $W_\tau^G$.
It is Hermitian and idempotent for a unitary realization, so its trace equals
its rank.  Therefore
\[
 \boxed{[m_\tau]\M_{G,V}=\dim(W_\tau^G).}
\]
This supplies both the invariant-multiplicity identity and the generalized
Molien formula used below.

\section{Representations of \texorpdfstring{$\Dic_m$}{Dic-m}}
Use
\[
 \Dic_m=\langle a,x:a^{2m}=1,\ x^2=a^m,\ xax^{-1}=a^{-1}\rangle,
 \qquad \zeta=e^{\pi i/m}.
\]
Its elements are $a^j$ and $xa^j$ for $0\le j<2m$.
The four one-dimensional representations are
\[
 L_{\alpha,\beta}:\quad a\mapsto\alpha,\quad x\mapsto\beta,
 \qquad \alpha^2=1,\quad\beta^2=\alpha^m.
\]
For $1\le k\le m-1$, the two-dimensional representations are
\[
 \rho_k(a)=\begin{pmatrix}\zeta^k&0\\0&\zeta^{-k}\end{pmatrix},\qquad
 \rho_k(x)=\begin{pmatrix}0&(-1)^k\\1&0\end{pmatrix}.
\]
The displayed matrices satisfy all three defining relations.  Their two
$a$-weights are distinct and are interchanged by $x$, hence $\rho_k$ is
irreducible.  The representations $\rho_k$, $1\leq k\leq m-1$, are pairwise
inequivalent: $\chi_k(a)=\zeta^k+\zeta^{-k}=2\cos(\pi k/m)$, and these values
are distinct because cosine is strictly decreasing on $(0,\pi)$.  They are
inequivalent to the one-dimensional representations by dimension.  The four
one-dimensional characters are pairwise distinct by their values on $a$ and
$x$.  Therefore the displayed irreducibles are pairwise inequivalent; now
$4\cdot1^2+(m-1)\cdot2^2=4m$ proves that the list is complete.

Write an arbitrary representation as
\[
 V\cong\bigoplus_{\alpha,\beta}L_{\alpha,\beta}^{\oplus c_{\alpha,\beta}}
 \oplus\bigoplus_{k=1}^{m-1}\rho_k^{\oplus d_k}.
\]
On $a^j$, the displayed summands have eigenvalues $\alpha^j$ and
$\zeta^{kj},\zeta^{-kj}$.  On $xa^j$, they have eigenvalue
$\beta\alpha^j$ and a pair whose square is $(-1)^k$.  Thus
\begin{align*}
\M_{\Dic_m,V}
={}&\frac1{4m}\sum_{j=0}^{2m-1}\prod_i
 \left[\prod_{\alpha,\beta}(1-\alpha^jt_i)^{-c_{\alpha,\beta}}
 \prod_k\frac1{(1-\zeta^{kj}t_i)^{d_k}(1-\zeta^{-kj}t_i)^{d_k}}\right]\\
&+\frac1{4m}\sum_{j=0}^{2m-1}\prod_i
 \left[\prod_{\alpha,\beta}(1-\beta\alpha^jt_i)^{-c_{\alpha,\beta}}
 \prod_k(1-(-1)^kt_i^2)^{-d_k}\right].
\end{align*}
This proves the arbitrary-representation formula directly from the spectrum.

\section{A single two-dimensional irreducible}
For $V=\rho_k$ the preceding expression becomes
\[
 \boxed{\M_{\Dic_m,\rho_k}=
 \frac1{4m}\sum_{j=0}^{2m-1}\prod_i
 \frac1{(1-\zeta^{kj}t_i)(1-\zeta^{-kj}t_i)}
 +\frac12\prod_i\frac1{1-(-1)^kt_i^2}.}
\]
Let $q_k=2m/\gcd(2m,k)$.  In the monomial basis
$u^{\tau_b-r_b}v^{r_b}$ of $\Sym^{\tau_b}V$, the $a$-weight is
$\zeta^{k(\tau_b-2r_b)}$.  Hence the cyclic-subgroup average selects exactly
\[
 N_{k,\tau}=\#\left\{(r_b):0\le r_b\le\tau_b,\quad
 q_k\mid |\tau|-2\sum_b r_b\right\}.
\]
For a coset element the eigenvalues are $q,-q$, where $q^2=(-1)^k$, and
\[
 h_d(q,-q)=\begin{cases}q^d,&d\text{ even},\\0,&d\text{ odd}.
 \end{cases}
\]
Half the group lies in each part, giving the exact integer formula
\[
 \boxed{[m_\tau]\M_{\Dic_m,\rho_k}
 =\frac12N_{k,\tau}+\frac12\mathbf1_{\text{all }\tau_b\text{ even}}
 (-1)^{k|\tau|/2}.}
\]
For $k=1$, the coset factor is $(1+t_i^2)^{-1}$, recording that lifted
reflections square to $-I$.

For $V=\bigoplus_k\rho_k^{\oplus d_k}$, the identical argument yields half a
zero-weight lattice count and half $\prod_bB_{\tau_b}$, where
\[
 B_d=[x^d]\prod_{k=1}^{m-1}(1-(-1)^kx^2)^{-d_k}.
\]

\section{Regular representation}
An element of order $d$ acts on the regular representation as $4m/d$ cycles
of length $d$.  There are $\varphi(d)$ elements of order $d$ in
$\langle a\rangle$, while every $xa^j$ has order four.  Therefore
\[
 \boxed{\M_{\Dic_m,\mathrm{reg}}=
 \frac1{4m}\sum_{d\mid2m}\varphi(d)
 \operatorname{Exp}_\sigma\!\left(\frac{4m}{d}p_d\right)
 +\frac12\operatorname{Exp}_\sigma(mp_4).}
\]

\section{Verification strategy and results}
The notebook uses an orthonormal occupancy basis for each symmetric power.
It embeds that basis in $V^{\otimes q}$, computes
$\Sym^q(\rho(g))$ by compression, forms the tensor product $W_\tau$, and then
forms the full Reynolds projector.  This route does not use eigenvalues or the
coefficient formula.  It is compared with (i) a character average, (ii) the
congruence count above, and (iii) coefficient extraction from the two spectral
sums.  A safety cap prevents accidental construction of oversized tensor
spaces.

\begin{center}
\small
\begin{tabular}{@{}lrrrr@{}}
\toprule
Case & $\dim W_\tau$ & rank & formula & error\\
\midrule
$\Dic_2,\rho_1,\tau=(2)$ & 3 & 0 & 0 & $6.8\cdot10^{-17}$\\
$\Dic_3,\rho_1,\tau=(2,2)$ & 9 & 2 & 2 & $6.9\cdot10^{-16}$\\
$\Dic_4,\rho_2,\tau=(3,1)$ & 8 & 2 & 2 & $3.4\cdot10^{-16}$\\
$\Dic_5,\rho_3,\tau=(4)$ & 5 & 1 & 1 & $5.5\cdot10^{-16}$\\
$\Dic_3,L_{1,-1}\oplus L_{-1,i}\oplus\rho_2,\tau=(2,1)$ & 40 & 4 & 4 & $<2.6\cdot10^{-15}$\\
$\Dic_4,\rho_1\oplus\rho_2,\tau=(2,1)$ & 40 & 2 & 2 & $<10^{-14}$\\
$\Dic_2,\mathrm{reg},\tau=(4)$ & 330 & 44 & 44 & $0$\\
\bottomrule
\end{tabular}
\end{center}
Here ``error'' is the projector idempotence residual or, for the two added
direct-sum rows, a conservative bound including relation and formula residuals.
The dimensions $2$, $4$, and $8$ exercise irreducible, reducible, and regular
representations without hiding the computation behind large numerical spaces.

\paragraph{Conclusion.}
Every tested route agrees.  More importantly, the proof identifies why the
formula must hold: the determinant series gives symmetric-power characters,
and averaging those characters is exactly the trace/rank of the invariant
projector.  The dicyclic sign is forced by $(xa^j)^2=a^m$.

\end{document}
